\section{Training pipeline}
\label{sec:training-pipeline}

Meowtion learns new behaviours through a closed loop that an owner drives entirely from the
dashboard, with no local toolchain: capture, label, train, deliver, run. The loop is what turns the
collar from a fixed-function sensor into a platform for recognising any behaviour the owner can
demonstrate.

\subsection{The loop}

\begin{enumerate}[leftmargin=1.6em,itemsep=2pt]
  \item \textbf{Define behaviours.} The owner lists the actions to recognise (eat, drink, purr,
        scratch, litter-tray, and so on) in the dev console. This set drives both the clip-label
        dropdown and the trained model's outputs.
  \item \textbf{Capture.} With a station in training mode, paired audio and time-aligned motion are
        recorded while a registered collar is in range and uploaded to Cloud Storage (Flow~2,
        \cref{sec:architecture}).
  \item \textbf{Label.} The owner reviews clips in the dev console, sets each clip's action, trims
        dead space from the waveform, and discards poor examples.
  \item \textbf{Train.} The \texttt{train} Cloud Function trains the models server-side from the
        labelled clips and writes the results back: the quantised \texttt{.tflite} blobs to Storage,
        and the new version and label list to the database.
  \item \textbf{Deliver.} Each station watches a model-version key; when it changes the station
        downloads the new \texttt{.tflite} files and pushes them to the collar over the air
        (\cref{sec:protocols}).
  \item \textbf{Run.} In production mode the collar classifies with the new model and the dashboard
        shows live, named behaviours.
\end{enumerate}

\subsection{The trainer}

Training runs in the \texttt{train} Python Cloud Function, gated by a developer account
(\cref{sec:cloud-backend}). It reads the account's labelled clips (metadata from the database, audio
and motion bytes from Storage), and derives the class set directly from the labels present, so the
network's outputs follow whatever behaviours the owner defined. It then trains the two cascade stages
independently:

\begin{itemize}[leftmargin=1.5em]
  \item a small 1-D convolutional network over the motion windows (six IMU channels), and
  \item a larger 1-D convolutional network over the 8\,kHz audio windows.
\end{itemize}

Each input is windowed with overlap and normalised exactly as the collar normalises it at inference,
so the on-device transform and the training transform match (\cref{sec:on-device-ai}). The trained
Keras models are converted to TensorFlow Lite with full int8 post-training quantisation using a
representative dataset, so the exported \texttt{.tflite} weights and activations are int8 and the
input and output tensors carry the scale and zero-point the collar uses to quantise its own inputs.

\subsection{Label-agnostic by design}

The trainer is not hard-coded to a fixed set of actions; it learns whatever label set the owner
defines. This is what makes Meowtion a platform rather than a fixed-function gadget: the same pipeline
that recognises eating and drinking can be pointed at any distinguishable behaviour the owner can
capture and label. The class set flows end to end as data, from the dashboard's action list, through
the trainer, to the index-based on-collar classifier, with names attached only at the edges.

\subsection{Delivery contract}

Training outputs are written to predictable locations so the station finds them without coordination:
the model blobs land in Storage as \texttt{models/<uid>/\allowbreak imu\_model.tflite} and
\texttt{models/<uid>/\allowbreak audio\_model.tflite}, the new version is mirrored into each of the owner's
stations' configuration so their version gate fires, and the label list is stored in
\texttt{users/<uid>/models} so the dashboard can map a class index back to a name across label-set
changes (\cref{sec:cloud-backend}).

\subsection{Known limitation}

A model trained on only two classes (for example, eat versus drink) has no way to represent
``neither'' and will confidently assign one of its two labels to an idle cat. The fix is to retrain
with an explicit idle/other class; this is a data and modelling task, not a firmware change, and is
tracked in \cref{sec:future-work}.
